\begin{abstract}
In September and November 2025 the self-propagating worms Shai-Hulud and Shai-Hulud~2.0 compromised
roughly \ValOneListed{} and \ValTwoListed{} npm packages. They did so by stealing maintainers' publish
tokens during package installation and republishing those maintainers' packages with an injected
install script. We model such worms as a two-channel SIR epidemic on the npm ecosystem.
\begin{itemize}
  \item \textbf{Dependency channel.} A compromised dependency exposes its dependents, weighted by how
        permissive each SemVer range is and by whether consumers block lifecycle scripts.
  \item \textbf{Credential channel.} A leaked token lets the worm republish up to a capped number of
        the owner's packages.
\end{itemize}

Our main data set is a \MainN-package graph built from the npm registry. It is the dependency closure
of the \NSeeds{} most-depended-upon packages, cleaned of a reward-farming spam ring.

\textbf{Calibration and validation.} With one tick equal to one hour, we calibrate transmission,
credential and patching rates to the observed plateau times and stolen-token share of both waves.
Against the real compromised-package lists, victims share a maintainer with another victim far more
often than random package sets do (\ValOneCoMaint{} vs.\ \ValOneCoMaintNull{} for wave~1, \ValTwoCoMaint{}
vs.\ \ValTwoCoMaintNull{} for wave~2). They are almost absent from the high-impact core.

\textbf{Findings.}
\begin{itemize}
  \item With dependencies alone, spread is bounded: outbreaks started at the most-depended-upon packages
        reach only \RQOneDepOnlyHub{} of the core, and leaf-started outbreaks fizzle.
  \item With the calibrated credential channel, even a leaf patient zero reaches \RQOneShOneLeaf{}
        (wave~1 parameters) to \RQOneShTwoLeaf{} (wave~2) of the core.
  \item Strict lockfiles and script blocking raise the dependency-channel epidemic threshold
        \RQTwoLockFactor$\times$ and \RQTwoScriptsFactor$\times$. With stealable tokens, however, the
        threshold in the dependency transmission rate collapses to the bottom of our grid.
  \item Reactive patching prevents almost all spread if the worm is detected within an hour, but only
        \RQThreeHubPatchingTwentyFour{} after a realistic 24-hour detection lag. Proactively immunizing
        the same number of high-degree packages still prevents \RQThreeHubImmunizationTwentyFour.
\end{itemize}
Credential hardening and proactive protection of central packages therefore matter more than dependency
pinning.
\end{abstract}
